\section{Classical Problems in Distributed Computing}
\label{app:additional-applications}

This appendix applies Lifted Monotonicity to classical problems of
distributed computing.  For each problem we give an \spectype: its inputs and
outputs, their orders, and its table.  We then find where Lifted Monotonicity
holds or fails, usually through the common-extension condition
(Proposition~\ref{prop:common-expansion-form}), and read off what Complete
CALM and Complete CAP say about it.  Classical results about failures,
computability, and termination depend on a model of computation, and none of
them follows from these specs.  We cite them where they bear on an example.

\subsection{Three patterns}
\label{app:classical-patterns}

Many classical specs follow one of three patterns.  In the first, outputs
report facts that only accumulate, and two reports permitted at one input can
be merged.  In the second, an output orders two concurrent events, and the
spec permits either order.  In the third, outputs decide values, and the spec
bounds how many distinct values may be decided.  The first pattern gives
Lifted Monotonicity, and each of the other two gives a witness to its failure.

\begin{lemma}[Accumulating outputs]
\label{lem:accumulating}
Suppose the outputs of an \spectype are sets of facts ordered by inclusion.  If
every output permitted at an input is still permitted at each \indesc of that
input, and the union of two outputs permitted at one input is permitted there,
then the spec is monotone.
\end{lemma}

\begin{proof}
Take $i\InputOrd j$ and a nonempty finite $X\subseteq\fiber{i}$.  Each member
of $X$ is permitted at $j$, so by induction on the size of $X$ their union is
permitted at $j$, and it \outrel every member of $X$.
\end{proof}

The proof uses only that each fiber contains, with any two outputs, a common
\outnoun of both.  The lemma therefore holds for any output order and any
binary operation with that property, such as the maximum of two counter
values under $\leq$.

Two events are \emph{concurrent} in an input when neither is causally before
the other, so the input does not order them.  If the spec lets an output order
them either way, two runs can choose opposite orders, and no later output
reconciles the two.

\begin{lemma}[Ordering concurrent events]
\label{lem:order-commitment}
Suppose that each output $o$ of an \spectype determines, for each pair of
events $a$ and $b$, whether $o$ \emph{places $a$ before $b$}; that if
$o\OutputOrd u$ and $o$ places $a$ before $b$, so does $u$; and that no output
permitted at an input $j$ places two events in both orders.  If rows
$(i_1,o_1)$ and $(i_2,o_2)$ have $j$ as a common \indesc, $o_1$ places $a$
before $b$, and $o_2$ places $b$ before $a$, then the spec is not monotone.  If
$i_1=i_2=j$, the fiber at $j$ is not directed.
\end{lemma}

\begin{proof}
A common \outnoun of $o_1$ and $o_2$ in $\fiber{j}$ would place $a$ before $b$
and $b$ before $a$.  So the common-extension condition fails, and by
Proposition~\ref{prop:common-expansion-form} so does Lifted Monotonicity.  When
$i_1=i_2=j$, the same argument shows that $o_1$ and $o_2$ are incompatible in
$\fiber{j}$.
\end{proof}

\begin{lemma}[Bounded number of values]
\label{lem:bounded-values}
Suppose that each output $o$ of an \spectype determines a set of values that
$o$ \emph{decides}; that if $o\OutputOrd u$, then $u$ decides every value that
$o$ decides; and that every output permitted at an input $j$ decides at most
$k$ values.  If $k+1$ rows have $j$ as a common \indesc and decide $k+1$
distinct values, one per row, then the spec is not monotone.
\end{lemma}

\begin{proof}
A common \outnoun of the $k+1$ outputs in $\fiber{j}$ would decide all $k+1$
values.  So the common-extension condition fails, and
Proposition~\ref{prop:common-expansion-form} applies.
\end{proof}

When Lemma~\ref{lem:order-commitment} or Lemma~\ref{lem:bounded-values}
applies, Complete CALM (Theorem~\ref{thm:complete-calm}) says that the runs of
the spec are not consistent.  The No-\commitment test
(Proposition~\ref{prop:exposure-boundary}) adds that some output, at some
state reachable without \commitment, is not legal, so emitting it applies a
\commitment.  With $E=P$, Theorem~\ref{thm:complete-cap} and
Corollary~\ref{cor:complete-cap-any-system} add that no system for the spec
has Availability for $P$ and Partition Consistency.  The witness rows, each
placed on its own side of a partition, show why.  When
Lemma~\ref{lem:accumulating} applies, the runs are consistent, and the
canonical system for $P$ needs no \commitment and escapes CAP for $P$.

Table~\ref{tab:classical-problems} classifies classical problems by pattern.
Snapshots, renaming, consensus, $k$-set agreement, and a causal FIFO queue are
treated in the subsections below.  We define the remaining specs here.

\begin{table*}[t]
  \centering
  \small
  \begin{tabular}{@{}lllll@{}}
    \toprule
    Problem & Input (order) & Output (order) & Lemma & Monotone \\
    \midrule
    Reliable broadcast & broadcast messages ($\subseteq$) & delivered messages ($\subseteq$) & \ref{lem:accumulating} & yes \\
    Causal broadcast, global snapshot & causal history & downward-closed set of events ($\subseteq$) & \ref{lem:accumulating} & yes \\
    Grow-only counter & increments ($\subseteq$) & count ($\leq$) & \ref{lem:accumulating}, with maximum & yes \\
    Total-order broadcast & broadcast messages ($\subseteq$) & delivery sequence (prefix) & \ref{lem:order-commitment} & no \\
    Causal FIFO queue & causal history of enqueues & dequeue sequence (prefix) & \ref{lem:order-commitment} & no; yes on chains \\
    Atomic snapshot & updates and scans ($\subseteq$) & views of answered scans (extension) & \ref{lem:order-commitment} & no \\
    Consensus & proposals ($\subseteq$) & decision map (extension) & \ref{lem:bounded-values}, $k=1$ & no \\
    $k$-set agreement & proposals ($\subseteq$) & decision map (extension) & \ref{lem:bounded-values} & no \\
    Leader election & candidates ($\subseteq$) & choice of leader (extension) & \ref{lem:bounded-values}, $k=1$ & no \\
    Renaming & participants ($\subseteq$) & names (extension) & \ref{lem:bounded-values}, $k=1$ & no \\
    \bottomrule
  \end{tabular}
  \caption{Classical problems by pattern.  Each row names the spec's inputs and
  outputs with their orders, and the lemma of
  Section~\ref{app:classical-patterns} that decides whether the spec is
  monotone.  A causal history is ordered by the downward-closed sub-history
  relation (Section~\ref{app:snapshots}), and ``chains'' are histories whose
  events are causally totally ordered (Section~\ref{app:queue}).}
  \label{tab:classical-problems}
\end{table*}

In \emph{reliable broadcast}, an input is a finite set of broadcast messages,
and an output is a set of delivered messages.  An output is permitted when it
contains only broadcast messages, so Lemma~\ref{lem:accumulating} applies.
Gossip and anti-entropy protocols~\cite{demers1987epidemic} propagate sets of
updates in the same way.
Reliable broadcast also requires that a message delivered by one correct
process is eventually delivered by every correct
process~\cite{hadzilacos1994modular}.  That requirement concerns termination,
and the spec does not express it.

A \emph{grow-only counter}~\cite{shapiro2011crdt} takes as input a finite set
of increments and outputs a count, ordered by $\leq$.  A count is permitted
when it is at most the number of increments.  A permitted count stays
permitted as increments arrive, and the larger of two permitted counts is
permitted, so the maximum form of Lemma~\ref{lem:accumulating} applies.

\emph{Total-order broadcast}~\cite{defago2004total} has the inputs of reliable
broadcast, but an output is a delivery sequence of distinct broadcast
messages, ordered by prefix.  Two delivery sequences have a common \outnoun
only when one is a prefix of the other, which is the agreement that
total-order broadcast requires.  A sequence places $a$ before $b$ when it lists
$a$, and lists $b$, if at all, after $a$.  An extension of the sequence keeps
this placement, and no sequence places two messages in both orders.  The rows
$(\{a\},[a])$ and $(\{b\},[b])$ have the common \indesc $\{a,b\}$, so by
Lemma~\ref{lem:order-commitment} the spec is not monotone: two processes that
each deliver their own message first cannot be reconciled.  In asynchronous
systems with crash failures, total-order broadcast and consensus are
equivalent problems~\cite{chandra1996unreliable}.

In \emph{leader election}, an input is a finite set of candidates, and an
output maps candidates to the leader each has chosen, ordered by extension.  An
output is permitted when it names only candidates and names at most one
leader.  This is consensus in which each candidate proposes itself, so
Lemma~\ref{lem:bounded-values} applies with $k=1$, as in
Section~\ref{app:consensus}.

\subsection{Consensus: committed authority, then monotone evidence}
\label{app:consensus}

As a spec, consensus is $k$-set agreement with $k=1$
(Section~\ref{app:kset}): inputs are proposals, outputs are decision maps, and
every permitted output decides at most one value.  Two participants that
propose different values and each decide their own give two rows with a
common \indesc, so by Lemma~\ref{lem:bounded-values} the spec is not monotone.
The impossibility result of Fischer, Lynch, and
Paterson~\cite{fischer1985impossibility} is a different statement: no
deterministic asynchronous protocol solves consensus if one process may crash.
That result concerns termination, which the spec does not express.

Practical consensus protocols such as Paxos~\cite{lamport2001paxos} and
Raft~\cite{ongaro2014raft} still make this \commitment, and their structure
determines who makes it.  Each epoch, a ballot in Paxos or a term in Raft, has an authority,
such as a leader, and a rule fixes which epoch succeeds which.  Within an
epoch, the authority chooses the value proposed for each slot of a replicated
log, and voters certify it.  Once the authority and the succession rule are
fixed, what remains to report is evidence: certified epochs, votes, and
certified decisions.  We show that reporting this evidence needs no further
\commitment.

Inputs are finite sets of protocol messages (proposals, votes,
certificates), ordered by inclusion.  An output is a triple $(E,V,D)$: a
certified epoch chain $E$, which is a finite sequence of epochs, each
certified under the authority of its predecessor; a set $V$ of votes or
certificates; and a set $D$ of decisions, at most one value per epoch and slot,
each certified for that epoch and slot.  Outputs are ordered by
\begin{equation}
  (E,V,D)\OutputOrd(E',V',D')
  \quad\Longleftrightarrow\quad
  E\preceq_{\rm pref}E',\quad V\subseteq V',\quad D\subseteq D',
  \label{eq:authority-refinement}
\end{equation}
where $E\preceq_{\rm pref}E'$ means that $E$ is a prefix of $E'$.  Write $c(i)$
for the certified evidence at input $i$: the epoch chain that the messages of
$i$ certify under the succession rule, the votes and certificates in $i$, and
the decisions certified in $i$.  The \emph{evidence spec} permits at $i$
exactly the outputs that $c(i)$ \outrel.  Protocol safety is the assumption
that $c(i)$ is defined at every input and grows with the input:
$i\InputOrd j$ implies $c(i)\OutputOrd c(j)$.  For Paxos, this is the
property that quorum intersection provides: once a value is certified for a
slot, no later certificate for that slot names another
value~\cite{lamport2001paxos}.

\begin{proposition}[Monotone evidence]
\label{prop:monotone-evidence}
If $i\InputOrd j$ implies $c(i)\OutputOrd c(j)$, the evidence spec is
monotone, and every output at every state reachable without \commitment is
legal.
\end{proposition}

\begin{proof}
Take rows $(i_r,o_r)$ with a common \indesc $j$.  Each
$o_r\OutputOrd c(i_r)\OutputOrd c(j)$, and $c(j)$ is permitted at $j$.  So the
common-extension condition holds, and
Proposition~\ref{prop:common-expansion-form} gives Lifted Monotonicity.  The
second claim follows by Proposition~\ref{prop:exposure-boundary}.
\end{proof}

The same reasoning applies at a single state.  If every row remaining at a
state \outrel an output $o$ of the current input, emitting $o$ removes no row,
so $o$ is a legal output.  Under protocol safety, this is the case for an
authority change or a certified decision once it is part of the evidence, so
emitting it needs no further \commitment.

The \commitments lie in the choices that the evidence spec takes as given.
Competing proposals for an undecided slot have incompatible extensions of $D$,
since a permitted output decides at most one value per epoch and slot.  So
certifying one of them is a \commitment.  If the succession rule were not
fixed, choosing among competing next epochs would require a \commitment as
well.  Protocols place these \commitments differently.  In Multi-Paxos and
Raft, the leader of the current epoch chooses the value for each slot.
Mencius~\cite{mao2008mencius} instead assigns each slot to a fixed replica in
advance, and that replica chooses the slot's value, possibly a no-op.

\subsection{$k$-set agreement: no $k$-wise test suffices}
\label{app:kset}

$k$-set agreement~\cite{chaudhuri1993more} weakens consensus: participants may
decide up to $k$ distinct values.  Weakening agreement in this way does not
restore monotonicity, and no test that examines $k$ proposal rows at a time
detects the failure.

Fix $k\geq1$, distinct participants $p_1,\ldots,p_{k+1}$ in the participant
set $\mathcal N$ (possibly among others), and distinct values
$v_1,\ldots,v_{k+1}$ in a value set $\mathit{Val}$.  An input is a finite set
of proposals $(p,v)$, at most one per participant, ordered by inclusion.  The
\emph{proposal input} of $p_r$ is $\{(p_r,v_r)\}$.  Outputs are partial
decision maps $m:\mathcal N\rightharpoonup\mathit{Val}$ ordered by extension.
At input $I$, $m$ is permitted when three conditions hold: it decides only
participants that propose in $I$, every value in its range is proposed in $I$
(validity), and its range has at most $k$ values.  At $p_r$'s proposal input,
$\{p_r\mapsto v_r\}$ is permitted, and the union of the proposal inputs is a
common \indesc.  Any $k$ of these $k+1$ rows have a common \outnoun there.  All
$k+1$ do not, since such an output would have $k+1$ values.  So the
common-extension condition fails, and the spec is not monotone
(Proposition~\ref{prop:common-expansion-form}).  This is
Lemma~\ref{lem:bounded-values}, where a map decides the values in its range.
A partition separating the $p_r$ gives the negative side of Complete CAP.
Consensus is $k=1$.

The failure lies within one fiber.  An output permitted at an input is still
permitted at each \indesc of that input, so Relational Monotonicity holds.  At
the union of the proposal inputs, each $\{p_r\mapsto v_r\}$ is permitted, but
the $k+1$ of them are not compatible in that fiber, so the fiber is not
directed.
Classical impossibility results for $k$-set agreement additionally depend on
the computation, failure, and progress model: it has no wait-free solution
when $k$ is less than the number of
participants~\cite{saks2000set,herlihy1999topological}.

\subsection{Renaming}
\label{app:renaming}

In renaming~\cite{attiya1990renaming}, participants with identifiers from a
large space choose distinct names from a small one.  In tight renaming the
namespace has exactly as many names as there are participants.  The
uniqueness of names, more than the size of the namespace, is what makes the
spec fail Lifted Monotonicity.

Inputs are finite sets $N$ of participants, ordered by inclusion.  Outputs are
finite partial naming maps ordered by extension, and the outputs at $N$ are
the injective maps
\begin{equation}
  m:N\longrightarrow\{1,\ldots,|N|\}.
  \label{eq:tight-renaming}
\end{equation}
Relational Monotonicity holds.  Given $N\subseteq N'$ and an output $m$ at $N$,
the participants of $N'\setminus N$ can take the names $|N|+1,\ldots,|N'|$,
which gives an output at $N'$ that \outrel $m$.

At $\{p\}$ the only output is $p\mapsto1$, and at $\{q\}$ it is $q\mapsto1$.
Each has an \outnoun at $\{p,q\}$, but no injective map \outrel both.  This is
Lemma~\ref{lem:bounded-values} with $k=1$, where a map decides the
participants to which it gives the name $1$.  The fiber at $\{p,q\}$ is not
directed either: $\{p\mapsto1,q\mapsto2\}$ and $\{p\mapsto2,q\mapsto1\}$ have
no common \outnoun.  The spec therefore fails Lifted Monotonicity, and by
Theorem~\ref{thm:complete-cap} with $E=P$ and
Corollary~\ref{cor:complete-cap-any-system} a partition separating $p$ and $q$
rules out Availability for $P$ with Partition Consistency.

The witness uses only the name $1$ at single-participant inputs and the
injectivity of outputs at $\{p,q\}$.  It therefore applies to renaming with any
namespace $\{1,\ldots,f(|N|)\}$ such that $f(n)\geq n$ for all $n$, for example
$f(n)=2n-1$.  The renaming literature separates such bounds by whether they can
be met wait-free~\cite{herlihy1991waitfree,herlihy1999topological}.  Those are
results about computation with crash failures, and the spec failure here does
not distinguish the bounds.

\subsection{Snapshots}
\label{app:snapshots}

Snapshots separate partial-order structure from total-order structure.  A
downward-closed cut of a causal history remains a cut as the history grows.
An atomic snapshot must, in contrast, order each scan relative to concurrent
updates.

Inputs are finite causal histories, that is, finite sets of environment events
with a causal partial order; $h\InputOrd h'$ when $h$ is a downward-closed
subset of $h'$ carrying the restricted order.  The \emph{prefix} spec permits
at $h$ any downward-closed subset of $h$, with outputs ordered by inclusion.
Old prefixes stay legal as history grows, and finite unions of prefixes are
prefixes, so it satisfies Lifted Monotonicity, by
Lemma~\ref{lem:accumulating}.  It guarantees neither freshness nor
completeness.  A global snapshot~\cite{chandy1985distributed} records such a
cut.  Causal broadcast delivers a message only after every message causally
before it, so its set of delivered messages is such a cut as well.

A spec that requires the complete history, with output $h$ at $h$ and outputs
ordered by equality, is functional and not monotone: its required output
changes whenever an environment event arrives.  By
Proposition~\ref{prop:functional-exposure-boundary}, some output can then be
emitted only after a \commitment that rules out larger inputs.

An atomic snapshot~\cite{afek1993atomic} answers scans with views, and the
answers must be explained by a single order of all scans and updates.  We
model one consequence of this requirement: any two views are comparable by
inclusion.  An input is a finite set of update and scan operations, ordered by
inclusion.  An output maps answered scans to views, each a set of updates of
the input, and is ordered by extension.  It is permitted when every answered
scan is in the input and any two of its views are comparable.  An output
places update $a$ before update $b$ when one of its views contains $a$ and not
$b$.  Let an input contain updates $u_1,u_2$ and scans $s_1,s_2$.  The outputs
$\{s_1\mapsto\{u_1\}\}$ and $\{s_2\mapsto\{u_2\}\}$ are both permitted there,
and they place $u_1$ and $u_2$ in opposite orders.  By
Lemma~\ref{lem:order-commitment}, the fiber is not directed and the spec is
not monotone.  This does not make atomic snapshots impossible: they have
wait-free implementations from atomic registers~\cite{afek1993atomic}.  Those
registers are linearizable, and the register example of
Section~\ref{sec:complete-cap} shows that a linearizable register fails
Lifted Monotonicity for its explicit relation.

\subsection{A causal FIFO queue}
\label{app:queue}

A replicated FIFO queue shows both outcomes within a single data type.
Enqueues that are causally ordered must be dequeued in that order.  Concurrent
enqueues are not ordered until some output orders them.

An input is a finite causal history of enqueue events, ordered as in
Section~\ref{app:snapshots}.  Each event enqueues a distinct value, which we
identify with the event.  An output is a dequeue sequence, ordered by prefix.
A sequence is permitted at $h$ when it lists distinct events of $h$, never
lists an event after one that is causally later, and skips no event: if it
lists an event, it lists every event causally before it.  We call this spec
the \emph{causal FIFO queue}.

\begin{proposition}[Causal FIFO queue]
\label{prop:causal-fifo}
The causal FIFO queue satisfies Relational Monotonicity.  Its fiber at a
history $h$ is directed if and only if any two events of $h$ are causally
ordered.  Hence the spec is not monotone if two values can be enqueued
concurrently, and its restriction to histories whose events are causally
totally ordered is monotone.
\end{proposition}

\begin{proof}
A sequence permitted at $h$ is still permitted at each $h'$ with
$h\InputOrd h'$.  The events of $h'$ outside $h$ are never causally before an
event of $h$, and the causal order on $h$ is unchanged.  So Relational
Monotonicity holds.

First assume that two events $a$ and $b$ of $h$ are concurrent.  The events
causally before $a$ can be listed in an order that respects causality, and
appending $a$ gives a permitted sequence that lists $a$ and not $b$.  In the
sense of Lemma~\ref{lem:order-commitment} it places $a$ before $b$.  A
permitted sequence that places $b$ before $a$ arises in the same way, so by
that lemma the fiber at $h$ is not directed.  Now assume that the events of $h$
are causally totally ordered.  In a permitted sequence, the event at each
position is the unique event whose causal predecessors are exactly the events
at earlier positions.  So two permitted sequences agree up to the length of
the shorter, which is therefore a prefix of the longer.  Any finite set of
permitted sequences then has its longest member as a common \outnoun, and the
fiber is directed.  The restriction to such histories has Relational
Monotonicity and directed fibers, which together are Lifted Monotonicity.
\end{proof}

At a history with concurrent enqueues $a$ and $b$, consider the two sequences
built in the proof.  Emitting the one that lists $a$ removes the row of the
one that lists $b$, and no remaining row at that input \outrel the removed
one.  So the output event
applies a \commitment, as ordering two transactions does in two-phase locking
(Section~\ref{sec:systems}).

Two variants show how the result depends on the spec.  A queue that reports
only the set of dequeued events, which must be downward closed, has the prefix
spec of Section~\ref{app:snapshots}, which is monotone: it does not report an
order among concurrent enqueues.  A queue without the no-skipping rule may
dequeue an event and never dequeue an event causally before it.  Its fiber at
any history with two distinct events $a$ and $b$ is not directed, even when $a$
is causally before $b$, since $[a]$ and $[b]$ are both permitted and place the
two events in opposite orders.

\begin{acks}
Thanks to Conor Power, whose persistent questioning of the CALM canon prompted this work.
Much gratitude to Peter Alvaro, Natacha Crooks, Chris Douglas, Tyler Hou, Martin
Kleppmann, Paris Koutris, Edward Lee, Shulu Li, Vincent Liu, David Chu McElroy,
Sarah Morin, and Dan Suciu for helpful feedback. Thanks also to
anonymous reviewers of this and previous submissions for frank feedback that pushed
this work deeper.

Generative AI tools (Infinity, Kiro, Claude and ChatGPT) were used as interactive writing and
reviewing assistants during the preparation of this work, including drafting
and revising text, and synthesizing Lean proofs for all mathematical arguments.  All
content was reviewed, validated, and approved by the author.
\end{acks}
